\section{Zusammenfassung und Ausblick}

\begin{frame}[plain]
    % \frametitle{<title>}
    \centering
    \Large
    \textbf{Zusammenfassung und Ausblick}
\end{frame}

\begin{frame}
    \frametitle{Zusammenfassung}
    \begin{itemize}
        \item Analytische Herleitung Selbstkonsistenzgleichung von $\Delta$
        \item Abhängigkeit $\Delta$ von Temperatur $T$, Debye-Energie $E_D$, Wechselwirkung $U$ und chemischem Potential $\mu$ 
        \begin{itemize}
            \item[$\rightarrow$] Bestätigung einer kritischen Wechselwirkung $U_C$ bei halber Füllung
            \item[$\rightarrow$] Dotierung zu Van-Hove-Singularitäten besonders interessant für Supraleitung
        \end{itemize}
        \item Anwendung auf Graphen
        \begin{itemize}
            \item[$\rightarrow$] keine konventionelle Supraleitung in reinem Graphen ($U_C \approx \qty{198}{eV}$)
            \item[$\rightarrow$] mögliche Supraleitung mit $T_C$ in Größenordnung $\num{0,1} - \qty{1}{K}$ an VHS
            \item[$\rightarrow$] Bandstruktur an VHS allgemein und unter starker Dotierung nicht gut beschrieben mit Modell
            \item[$\rightarrow$] Bestimmung von $U$ schwierig 
            \item[$\implies$] Aussagekraft beeinträchtigt
        \end{itemize}
    \end{itemize}
\end{frame}

\begin{frame}
    \frametitle{Ausblick}

    \begin{itemize}
        \item Anwendung auf weitere Materialien ($\mathrm{MoS_2}$, Bornitrid, ...)
        \begin{itemize}
            \item[$\rightarrow$] Modell muss dafür erweitert werden (höhere Tight-Binding-Terme, Berücksichtigung von mehreren Orbitalen)
        \end{itemize}
        \item Erweiterung auf nicht-konventionelle Supraleitung ($d$-wave, etc.)
        % \begin{itemize}
        %     \item[$\rightarrow$] Hinweise dazu auch in Graphen
        % \end{itemize}
    \end{itemize}

\end{frame}